% ============================================================================
%   APPENDIX B — FULL PROMPTS
% ============================================================================

\chapter{Full Prompt Specifications}
\label{app:prompts}

This appendix reproduces verbatim the four prompt variants used in the 
study, in the exact form applied to the models during inference. Two 
prompts are described (Prompt~A and Prompt~B), each with two 
formulations depending on the target model class 
(instruction-tuned or base). The dynamic fields \texttt{\{note\}}, 
\texttt{\{criterion\}}, and \texttt{\{criterion\_type\}} are substituted 
at query time with the content of each patient--criterion pair. All 
patient notes are truncated to $1{,}500$ characters before insertion to 
control for context-window variability across models.

% ----------------------------------------------------------------------------
\section{Prompt A --- Generic Baseline}
\label{sec:app-prompt-a}

Prompt~A is a minimalist prompt containing only the classification 
instruction, the input note, the criterion text, and the list of 
candidate labels. It does not define the labels, does not provide 
examples, and does not include reasoning rules.

\subsection{Prompt A for Instruct Models}

For instruction-tuned models, Prompt~A is wrapped in the model's native 
chat template via \texttt{tokenizer.apply\_chat\_template} with the 
generation prompt enabled.

\begin{lstlisting}[caption={Prompt A for Instruct models (user 
    message before chat-template wrapping)}, 
    label={lst:prompt-a-instruct}, breaklines=true, basicstyle=\ttfamily\small]
Classify whether this patient meets the eligibility criterion.

Patient: {note}
Criterion ({criterion_type}): {criterion}

Choose exactly one label from: included, not included, excluded,
not excluded, not enough information, not applicable.

Answer:
\end{lstlisting}

\subsection{Prompt A for Base Models}

For base models, no chat template is applied. Instead, a single 
few-shot example is embedded in the prompt to establish the desired 
input--output pattern by demonstration.

\begin{lstlisting}[caption={Prompt A for Base models}, 
    label={lst:prompt-a-base}, breaklines=true, basicstyle=\ttfamily\small]
Task: Classify whether a patient meets one eligibility criterion.
Choose exactly one label from: included, not included, excluded,
not excluded, not enough information, not applicable.

Example:
Patient note: 65-year-old male with type 2 diabetes.
Criterion (inclusion): diabetes
Final label: included

Now classify:
Patient note: {note}
Criterion ({criterion_type}): {criterion}
Final label:
\end{lstlisting}

% ----------------------------------------------------------------------------
\section{Prompt B --- Adapted with Expert Guidance}
\label{sec:app-prompt-b}

Prompt~B is an adapted prompt containing definitions of the labels, 
one worked example per label, and a critical inference rule for 
exclusion criteria. The prompt varies structurally between inclusion 
and exclusion criteria: two dedicated guidance blocks (one per 
criterion type) are inserted into a common frame. This design 
addresses the tendency of language models to over-use ``not enough 
information'' for exclusion criteria whose target condition is not 
explicitly discussed in the patient note.

\subsection{Prompt B for Instruct Models}

For instruction-tuned models, Prompt~B is again wrapped in the 
model's native chat template. The full prompt is constructed by 
concatenating the criterion-type-specific guidance block with the 
common frame shown below.

\begin{lstlisting}[caption={Prompt B for Instruct models --- common 
    frame}, label={lst:prompt-b-instruct-frame}, 
    breaklines=true, basicstyle=\ttfamily\small]
You are an expert clinical trial coordinator. Classify whether a
patient meets one eligibility criterion, based only on the patient
note.

{guidance block for inclusion or exclusion, see below}

KEY DISTINCTION:
- Use "not enough information" ONLY when the note is silent and you
  cannot reasonably assume the answer.
- If the note gives ANY evidence (even indirect) about the patient's
  status, use a decisive label (included/not included or
  excluded/not excluded).

---
PATIENT NOTE:
{note}

CRITERION ({criterion_type}):
{criterion}
---

Reason about whether the note addresses this criterion, then give
your final label.
Final label:
\end{lstlisting}

\paragraph{Guidance block for inclusion criteria.}
The following block replaces the placeholder in the common frame when 
the criterion type is \emph{inclusion}.

\begin{lstlisting}[caption={Prompt B for Instruct models --- inclusion 
    guidance}, label={lst:prompt-b-instruct-incl}, 
    breaklines=true, basicstyle=\ttfamily\small]
You are evaluating an INCLUSION criterion. To join the trial, the
patient SHOULD satisfy it.

Choose exactly ONE label:

- "included" - The patient note contains information showing the
  patient MEETS this inclusion criterion.
  Example: criterion "Type 2 diabetes", note says "patient has T2DM"
  -> included.

- "not included" - The patient note contains information showing the
  patient does NOT meet this inclusion criterion.
  Example: criterion "age over 65", note says "45-year-old patient"
  -> not included.

- "not enough information" - The note does NOT mention anything that
  lets you decide. The information is simply missing.
  Example: criterion "HbA1c between 7 and 10", note never mentions
  HbA1c -> not enough information.

- "not applicable" - The criterion concerns a different population
  or context and cannot logically apply to this patient.
  Example: criterion "for pediatric patients only", note describes
  an adult -> not applicable.
\end{lstlisting}

\paragraph{Guidance block for exclusion criteria.}
The following block replaces the placeholder in the common frame when 
the criterion type is \emph{exclusion}. This block contains the 
critical inference rule for exclusion criteria referenced in 
Section~\ref{subsec:met-prompt-b} of Chapter~\ref{chap:methodology}.

\begin{lstlisting}[caption={Prompt B for Instruct models --- exclusion 
    guidance}, label={lst:prompt-b-instruct-excl}, 
    breaklines=true, basicstyle=\ttfamily\small]
You are evaluating an EXCLUSION criterion. To join the trial, the
patient should NOT have the condition described.

Choose exactly ONE label:

- "excluded" - The note shows the patient HAS the excluding
  condition. They ARE excluded.
  Example: criterion "pregnancy", note says "patient is pregnant"
  -> excluded.

- "not excluded" - The patient does NOT have the excluding
  condition, so they can join (FAVORABLE).
  This includes TWO cases:
  (a) The note explicitly says the patient does not have it
      ("no history of cancer").
  (b) The note describes the patient's relevant conditions but does
      NOT mention this one - for a specific medical condition,
      absence of mention means the patient does not have it.
  Example: criterion "history of stroke", note details the
  patient's history with no stroke mentioned -> not excluded.

- "not enough information" - Use this RARELY, only when the
  criterion needs a specific measurement or detail that the note
  genuinely cannot address and that cannot be assumed absent.
  Example: criterion "ejection fraction below 40%", note never
  measures ejection fraction -> not enough information.

- "not applicable" - The criterion targets a different population
  and cannot logically apply.
  Example: criterion "pregnant women excluded", patient is male
  -> not applicable.

CRUCIAL RULE FOR EXCLUSION CRITERIA:
If the criterion names a specific disease, condition, or medical
history, and the note does NOT mention it, the default is
"not excluded" (the patient most likely does not have it). Reserve
"not enough information" only for criteria requiring a specific lab
value or measurement that is truly missing.
\end{lstlisting}

\subsection{Prompt B for Base Models}

For base models, the enriched-guidance structure of the instruct 
version is replaced with a set of few-shot examples. Three examples 
are provided for inclusion criteria, and four for exclusion criteria; 
the additional exclusion example encodes the inference rule for 
exclusion criteria by demonstration.

\begin{lstlisting}[caption={Prompt B for Base models --- common 
    frame}, label={lst:prompt-b-base-frame}, 
    breaklines=true, basicstyle=\ttfamily\small]
Task: Classify whether a patient meets one eligibility criterion,
based only on the patient note. Choose exactly one label from:
included, not included, excluded, not excluded, not enough
information, not applicable.

{few-shot block for inclusion or exclusion, see below}

Now classify this case:
Patient note: {note}
Criterion ({criterion_type}): {criterion}
Final label:
\end{lstlisting}

\paragraph{Few-shot block for inclusion criteria.}

\begin{lstlisting}[caption={Prompt B for Base models --- inclusion 
    few-shot examples}, label={lst:prompt-b-base-incl}, 
    breaklines=true, basicstyle=\ttfamily\small]
Example 1:
Patient note: 65-year-old male with type 2 diabetes, HbA1c 8.5%.
Criterion (inclusion): Type 2 diabetes mellitus
Final label: included

Example 2:
Patient note: 45-year-old female, no diabetes history.
Criterion (inclusion): age over 65
Final label: not included

Example 3:
Patient note: 50-year-old patient with hypertension.
Criterion (inclusion): HbA1c between 7 and 10
Final label: not enough information
\end{lstlisting}

\paragraph{Few-shot block for exclusion criteria.}
The fourth example (Example~3, ``history of stroke'') encodes the 
exclusion-inference rule of Prompt~B: when the criterion names a 
specific condition that the note does not mention, the model should 
default to ``not excluded'' rather than to ``not enough information''.

\begin{lstlisting}[caption={Prompt B for Base models --- exclusion 
    few-shot examples}, label={lst:prompt-b-base-excl}, 
    breaklines=true, basicstyle=\ttfamily\small]
Example 1:
Patient note: 30-year-old female, 12 weeks pregnant.
Criterion (exclusion): pregnancy
Final label: excluded

Example 2:
Patient note: 55-year-old male with hypertension and diabetes,
no cancer history.
Criterion (exclusion): history of cancer
Final label: not excluded

Example 3:
Patient note: 60-year-old male with hypertension. Otherwise healthy.
Criterion (exclusion): history of stroke
Final label: not excluded

Example 4:
Patient note: 70-year-old male with chronic kidney disease.
Criterion (exclusion): ejection fraction below 40%
Final label: not enough information
\end{lstlisting}